%==============================================================
\section{Ορισμός Προβλήματος}
%==============================================================

Η παρούσα ενότητα εξειδικεύει τον τυπικό μαθηματικό ορισμό του Πολύστροφου RAG και αναλύει τις δομικές προκλήσεις που αυτό συνεπάγεται. Το πρόβλημα προσεγγίζεται ως μια σύνθετη διαδικασία διαδοχικής λήψης αποφάσεων υπό συνθήκες στοχαστικής αβεβαιότητας, η οποία απαιτεί τη σύζευξη δυναμικής ανάκτησης πληροφορίας, παραγωγής φυσικής γλώσσας υπό περιορισμούς (constrained generation) και στοχευμένης εκτίμησης απαντησιμότητας (answerability estimation).

%--------------------------------------------------------------
\subsection{Τυπική Διατύπωση Πολύστροφου RAG}
%--------------------------------------------------------------

Έστω ένας διάλογος εξελισσόμενος σε διαδοχικές στροφές (turns) $i = 1, 2, \dots, t$. Ορίζουμε το διαλογικό ιστορικό $H_t$ μέχρι τη χρονική στιγμή $t$ ως την ακολουθία των προηγούμενων αλληλεπιδράσεων:

\begin{equation}
H_t = \Big( (u_1, \mathcal{D}_1, y_1), (u_2, \mathcal{D}_2, y_2), \dots, (u_{t-1}, \mathcal{D}_{t-1}, y_{t-1}) \Big)
\end{equation}

όπου $u_i$ αντιπροσωπεύει το εκφώνημα (utterance) του χρήστη, $\mathcal{D}_i \subset \mathcal{D}$ το σύνολο των εγγράφων που ανακτήθηκαν και $y_i$ την απόκριση της στροφής $i$. Στο MTRAGEval το ιστορικό αυτό είναι το ιστορικό αναφοράς (reference history) και όχι οι απαντήσεις του ίδιου του συστήματος. Στην τρέχουσα στροφή $t$, ο χρήστης εισάγει ένα νέο, ενδεχομένως μη αυτοτελές ερώτημα $q_t \equiv u_t$.

Με δεδομένο ένα στατικό, καθολικό σώμα γνώσης (knowledge corpus) $\mathcal{D}$, η αρχιτεκτονική Πολύστροφου RAG αποσυντίθεται σε τρεις διασυνδεδεμένες συναρτήσεις, οι οποίες αντιστοιχούν απευθείας στις προδιαγραφές του MTRAGEval:

\paragraph{1. Συνάρτηση Ανάκτησης Συμφραζομένων (Task A):}
Ο μηχανισμός ανάκτησης $\mathcal{R}$ αναλαμβάνει τη χαρτογράφηση του τρέχοντος ερωτήματος και του ιστορικού σε ένα βέλτιστο, διατεταγμένο υποσύνολο $k$ σχετικών αποσπασμάτων (κατά τις προδιαγραφές του διαγωνισμού, $k=10$):
\begin{equation}
\mathcal{R}(q_t, H_t, \mathcal{D}) \rightarrow \mathcal{D}_t = \{d_1, d_2, \dots, d_k\} \subset \mathcal{D}, \quad k=10
\end{equation}

\paragraph{2. Συνάρτηση Τεκμηριωμένης Παραγωγής (Task B):}
Με δεδομένο ένα χρυσό (ground-truth) σύνολο αναφοράς αποδείξεων $\mathcal{D}_t^*$, ο μηχανισμός παραγωγής $\mathcal{G}$ συνθέτει μια πραγματολογικά αγκυρωμένη απόκριση $y_t$, απομονώνοντας πλήρως την ποιότητα παραγωγής από τυχόν σφάλματα ανάκτησης:
\begin{equation}
\mathcal{G}(q_t, H_t, \mathcal{D}_t^*) \rightarrow y_t
\end{equation}

\paragraph{3. Ενοποιημένη End-to-End Παραγωγή (Task C):}
Στο πλήρες σύστημα, η παραγωγή εξαρτάται άμεσα από τα αποτελέσματα του πρώτου σταδίου, γεγονός που εισάγει στοχαστική εξάρτηση και επιτρέπει τη σώρευση σφαλμάτων ανάκτησης στην τελική απόκριση:
\begin{equation}
P(y_t \mid q_t, H_t) = \sum_{\mathcal{D}_t} P(y_t \mid q_t, H_t, \mathcal{D}_t) \cdot P(\mathcal{D}_t \mid q_t, H_t)
\end{equation}

%--------------------------------------------------------------
\subsection{Το Πρόβλημα της Απαντησιμότητας (Answerability Gate)}
%--------------------------------------------------------------

Σε πραγματικά σενάρια συνομιλίας, η πληροφοριακή ανάγκη του χρήστη ενδέχεται να μην καλύπτεται από το διαθέσιμο corpus $\mathcal{D}$. Για τη μοντελοποίηση αυτής της συνθήκης, εισάγουμε μια λανθάνουσα δυαδική μεταβλητή απαντησιμότητας $z_t \in \{0, 1\}$, η οποία ορίζεται ως:

\begin{equation}
z_t =
\begin{cases}
1, & \text{εάν το } \mathcal{D} \text{ περιέχει επαρκή στοιχεία για την πλήρη τεκμηρίωση του } q_t \text{ υπό το } H_t \\
0, & \text{διαφορετικά (unanswerable query)}
\end{cases}
\end{equation}

Το σύστημα οφείλει να εκτιμήσει το $z_t$ μέσω μιας δυαδικής απόφασης $\hat{z}_t \in \{0,1\}$. Στο προτεινόμενο σύστημα, κάθε κριτής $j$ αποφαίνεται για την απαντησιμότητα με βαθμό βεβαιότητας $c_j$, και η στροφή κρίνεται μη απαντήσιμη όταν κάποιος κριτής δηλώνει ανεπάρκεια στοιχείων με βεβαιότητα τουλάχιστον ίση με ένα κατώφλι βαθμονόμησης $\tau$:
\begin{equation}
\hat{z}_t =
\begin{cases}
0, & \text{εάν } \exists j: \text{ο κριτής } j \text{ δηλώνει ανεπάρκεια με } c_j \geq \tau \\
1, & \text{διαφορετικά}
\end{cases}
\end{equation}

Η τελική απόκριση προκύπτει ως:
\begin{equation}
y_t = 
\begin{cases}
\mathcal{G}(q_t, H_t, \mathcal{D}_t), & \text{εάν } \hat{z}_t = 1 \\
\text{«Δεν διαθέτω επαρκείς πληροφορίες για να απαντήσω»}, & \text{εάν } \hat{z}_t = 0
\end{cases}
\end{equation}

Το σφάλμα στην εκτίμηση του $z_t$ παρουσιάζει έντονη ασυμμετρία: τα ψευδώς θετικά αποτελέσματα ($\hat{z}_t=1$ ενώ $z_t=0$) αναγκάζουν το ΜΓΜ να υποπέσει σε πραγματολογικές ψευδαισθήσεις λόγω της τάσης συμμόρφωσης (compliance bias), ενώ τα ψευδώς αρνητικά μειώνουν δραστικά τη χρηστικότητα του συστήματος. Όπως δείχνουμε στην Ενότητα 1.5.4, η ασυμμετρία αυτή δεν είναι θεωρητική μόνο: στο πραγματικό σύστημα εκδηλώνεται ως ισχυρή μεροληψία αποδοχής (acceptance bias) προς το $\hat{z}_t=1$.

%--------------------------------------------------------------
\subsection{Διατύπωση ως Πρόβλημα Ενιαίας Βελτιστοποίησης}
%--------------------------------------------------------------

Η συνολική αρχιτεκτονική στοχεύει στη μεγιστοποίηση της αναμενόμενης χρησιμότητας της απόκρισης υπό μια πολυδιάστατη συνάρτηση κόστους:

\begin{equation}
\mathcal{O}^* = \arg\max_{\mathcal{R}, \mathcal{G}} \mathbb{E}_{(q_t, H_t) \sim \mathcal{X}} \Big[ \mathcal{U}\big(y_t, y_t^*, \mathcal{D}_t, \mathcal{D}_t^*\big) \Big]
\end{equation}

Η συνάρτηση χρησιμότητας $\mathcal{U}$ αναλύεται σε τέσσερις ανεξάρτητους πυλώνες:
\begin{equation}
\mathcal{U} = w_1 \cdot \text{NDCG}(\mathcal{D}_t, \mathcal{D}_t^*) + w_2 \cdot \text{Faithfulness}(y_t, \mathcal{D}_t) + w_3 \cdot \text{Extractive\_F1}(y_t, y_t^*) + w_4 \cdot \Phi(\hat{z}_t, z_t)
\end{equation}
όπου το $\Phi$ ανταμείβει τις σωστά βαθμονομημένες αποφάσεις της πύλης απαντησιμότητας (συμφωνία του $\hat{z}_t$ με το $z_t$)· αφού η $\mathcal{U}$ μεγιστοποιείται, όλοι οι όροι εκφράζονται ως ανταμοιβές. Τα βάρη $w_1, \dots, w_4$ δεν εκτιμώνται αριθμητικά: η διατύπωση είναι εννοιολογική. Λόγω της μη διαφορίσιμης φύσης των μετρικών κειμένου και της ανάκτησης, η βελτιστοποίηση αυτή δεν επιδιώκεται μέσω end-to-end gradient descent, αλλά προσεγγίζεται μέσω σύνθετων agentic pipelines και στοχευμένης μηχανικής προτροπών (prompt engineering), όπως περιγράφεται στην Ενότητα 1.4.

%--------------------------------------------------------------
\subsection{Ιδιαιτερότητες του Πολύστροφου Περιβάλλοντος}
%--------------------------------------------------------------

Το πολύστροφο RAG εισάγει τρεις θεμελιώδεις ακαδημαϊκές προκλήσεις:
\begin{itemize}
    \item \textbf{Χάσμα Λεξιλογίου (Vocabulary Gap):} Τα ερωτήματα μεταγενέστερων στροφών είναι σύντομα και ελλειπτικά. Η απουσία ρητών λέξεων-κλειδιών καθιστά τους κλασικούς ανακτητές (π.χ. BM25) ανίκανους να γεφυρώσουν τη σημασιολογική απόσταση με το corpus.
    \item \textbf{Σωρευτική Διάδοση Σφαλμάτων (Error Cascade):} Σε έναν πραγματικό βοηθό, εάν η συνάρτηση $\mathcal{R}$ αποτύχει στη στροφή $t-1$, ο θόρυβος ενσωματώνεται στο $H_t$ μέσω της απάντησης του συστήματος, και η επόμενη αναδιατύπωση βασίζεται σε μολυσμένα συμφραζόμενα. Το MTRAGEval παρέχει το ιστορικό αναφοράς σε κάθε στροφή, οπότε ο μηχανισμός αυτός δεν ενεργοποιείται στην αξιολόγησή μας· παραμένει όμως η δυσκολία επίλυσης αναφορών σε προηγούμενες στροφές.
    \item \textbf{Μετατόπιση Πλαισίου (Context Drift):} Η δυναμική εναλλαγή θεμάτων εντός του ίδιου διαλόγου απαιτεί από το σύστημα να απομονώνει το σχετικό ιστορικό και να αγνοεί τις ξεπερασμένες πληροφορίες.
\end{itemize}

%--------------------------------------------------------------
\subsection{Κεντρική Υπόθεση Εργασίας}
%--------------------------------------------------------------

Η παρούσα διπλωματική εργασία θεμελιώνεται πάνω στην εξής κεντρική ερευνητική υπόθεση:
\begin{quote}
\textit{«Η συνολική απόδοση ενός end-to-end συστήματος Πολύστροφου RAG φράσσεται άνω (upper-bounded) από τη σταθερότητα της διαδικασίας ανάκτησης (retrieval stability) και τη βαθμονόμηση της πύλης απαντησιμότητας, και όχι από την εγγενή εκφραστική ή συνθετική ικανότητα των σύγχρονων γλωσσικών μοντέλων.»}
\end{quote}
Συνεπώς, η έρευνα εστιάζει στη σταθεροποίηση του $P(\mathcal{D}_t \mid q_t, H_t)$ μέσω ελεγχόμενης ποικιλότητας ερωτημάτων, παρά στην αναζήτηση ισχυρότερων ή μεγαλύτερων γλωσσικών μοντέλων.

%--------------------------------------------------------------
\subsection{Συσχέτιση με το Benchmark MTRAGEval}
%--------------------------------------------------------------

Η αξιολόγηση των προτεινόμενων λύσεων πραγματοποιείται στο απαιτητικό περιβάλλον του SemEval-2026 Task 8, το οποίο αντανακλά τις ανωτέρω προκλήσεις σε ρεαλιστική κλίμακα: (α) η συντριπτική πλειονότητα των στροφών ($87\%$) είναι μη αυτοτελείς, με αναφορές αντωνυμίας να εμφανίζονται στο $97\%$ των μη-πρώτων στροφών· (β) το test set περιλαμβάνει σχεδόν τριπλάσιο ποσοστό μη απαντήσιμων στροφών σε σχέση με το development set· και (γ) παρατηρείται έντονη κατανομική μετατόπιση (distribution shift) μεταξύ development και test set τόσο ως προς την αναλογία unanswerable ερωτημάτων όσο και ως προς την κατανομή των θεματικών πεδίων (domains). Ο Πίνακας~\ref{tab:dev-test-greek} συνοψίζει τα κρίσιμα αυτά μεγέθη.

\begin{table}[h]
\centering
\small
\begin{tabular}{lcc}
\toprule
\textbf{Χαρακτηριστικό} & \textbf{Development} & \textbf{Test} \\
\midrule
Συνομιλίες & 110 & 507 \\
Σύνολο αξιολογούμενων στροφών & 842 & 507 \\
Μέσος αριθμός στροφών/συνομιλία & 7.7 & 1.0 \\
Μη αυτοτελείς στροφές & $87\%$ & — \\
Απαντήσιμες ή μερικώς απαντήσιμες & $92.3\%$ & $74.4\%$ \\
\textbf{Μη απαντήσιμες} & $6.5\%$ & $\mathbf{19.1\%}$ \\
Συνομιλιακές / λοιπές & $1.2\%$ & $6.5\%$ \\
\bottomrule
\end{tabular}
\caption{Σύγκριση development και test set στο MTRAGEval. Η σχεδόν τριπλάσια αύξηση του ποσοστού μη απαντήσιμων στροφών στο test set αποτελεί την κύρια πηγή της κατανομικής μετατόπισης που αναλύεται στην Ενότητα 1.5.5. Στο test set, $377$ στροφές διαθέτουν κρίσεις συνάφειας για το Task~A και $97$ είναι μη απαντήσιμες· η κατηγορία των υπόλοιπων $33$ στροφών ($6.5\%$) δεν μπόρεσε να ανακτηθεί από τα διαθέσιμα αρχεία.}
\label{tab:dev-test-greek}
\end{table}

%--------------------------------------------------------------
\subsection{Ερευνητικά Ερωτήματα}
%--------------------------------------------------------------

Με βάση τη διατύπωση του προβλήματος, η εργασία επιχειρεί να απαντήσει στα εξής τέσσερα ερευνητικά ερωτήματα (Research Questions):
\begin{enumerate}
    \item[\textbf{RQ1:}] Μπορεί η εισαγωγή στοχευμένης ποικιλότητας ερωτημάτων (query diversity) σε έναν ενιαίο ανακτητή να υποκαταστήσει ή να υπερβεί τη σύνθεση ετερογενών μοντέλων ανάκτησης (retriever diversity);
    \item[\textbf{RQ2:}] Σε ποιο βαθμό η αποσύνθεση της παραγωγής με εισαγωγή σταδίου εξαγωγής αποσπασμάτων (verbatim span extraction) συμβάλλει στον περιορισμό των πραγματολογικών ψευδαισθήσεων;
    \item[\textbf{RQ3:}] Ποια είναι η επίδραση της κατανομικής μετατόπισης των unanswerable ερωτημάτων στην ευαισθησία των πυλών απαντησιμότητας;
    \item[\textbf{RQ4:}] Μπορεί η επαναβαθμολόγηση μέσω ΜΓΜ (LLM-based reranking) να προσφέρει μετρήσιμο όφελος έναντι εξειδικευμένων cross-encoder μοντέλων ανάκτησης, όταν η ανάκληση του πρώτου σταδίου ανάκτησης πλησιάζει τον κορεσμό;
\end{enumerate}
